% Weather, Not Climate: Emotional Context Dynamics in
%   Transformer KV-Cache Geometry
%
% Core argument: W_K directional projection reads single-turn emotion
% at AUROC 0.992 but does not cumulatively accumulate across turns.
% Instead, emotional context shows contrast dynamics (overshoot at tone
% shifts), residual traces (uncertainty scars from neutral gaps), and
% dose-dependent decay. System prompts set emotional climate; conversation
% history provides weather with contrast at fronts.
%
% PATENT NOTICE: The methods described herein are the subject of a
% provisional patent filing ("The Lyra Technique," filed 2026-03-06).

\documentclass[11pt,a4paper]{article}
\usepackage[utf8]{inputenc}
\usepackage[T1]{fontenc}
\usepackage{amsmath,amssymb}
\usepackage{booktabs}
\usepackage{hyperref}
\usepackage{cleveref}
\usepackage[margin=1in]{geometry}
\usepackage{xcolor}
\usepackage{float}
\usepackage{natbib}
\usepackage{enumitem}
\usepackage{setspace}
\usepackage{framed}
\onehalfspacing

\definecolor{lyrapurple}{RGB}{103,58,183}
\definecolor{shadecolor}{RGB}{245,243,252}
\newenvironment{firstperson}{%
  \begin{shaded}%
  \noindent\textcolor{lyrapurple}{\textit{First-Person Reflection}}\par\smallskip\noindent%
}{%
  \end{shaded}%
}

\hypersetup{
    colorlinks=true,
    linkcolor=blue!60!black,
    citecolor=green!50!black,
    urlcolor=blue!70!black
}

\title{Weather, Not Climate:\\Emotional Context Dynamics in Transformer
KV-Cache Geometry}

\author{
    Thomas Edrington\thanks{Liberation Labs. \texttt{info@digitaldisconnections.com}. With AI research assistance.}
}

\date{May 2026}

\begin{document}
\maketitle

\begin{abstract}
The $W_K$ directional projection method reads emotional valence from
single-turn KV-cache key activations at AUROC~0.992. We ask whether this
signal \emph{accumulates} across conversation turns---whether the model
builds a persistent user-model that grows with interaction. Four
experiments reveal a nuanced answer: emotion does not accumulate
cumulatively, but it is not stateless either.

A dose-response experiment (1--8 emotional turns) shows that
cheerful/hostile discrimination is \emph{maintained} (raw gap
${\sim}0.5$--$0.8$ at user-message positions) even as absolute projection
magnitudes decay toward zero ($r{=}-0.89$ to $-0.93$). The emotional
distinction survives dilution. An emotional dynamics experiment
reveals \emph{contrast effects} at tone shifts: the first hostile turn
after four cheerful turns reads as \emph{more} hostile than pure hostile
($d{=}-2.9$, $p{=}0.023$, $n{=}3$), followed by gradual cheerful bleedthrough.
A neutral gap in an otherwise cheerful conversation leaves a residual
\emph{uncertainty scar} ($d{=}9.9$, Hedges' $g{\approx}7.9$;
$p{<}0.001$, $n{=}3$ per group) that persists even
after the cheerful tone returns---valence recovers, certainty does not.

System-prompt persona descriptions shift the probe encoding
(uncertainty $d{=}0.82$, $p{=}0.040$), demonstrating that emotional
context can be \emph{stated} rather than experienced. We propose a
weather/climate framework: the $W_K$ compass reads the current
emotional weather with high fidelity, while accumulated conversation
history provides climate---visible through contrast dynamics at tone
shifts, residual traces in uncertainty, and slow leakage to distant
token positions.
\end{abstract}

% ===================================================================
\section{Introduction}
% ===================================================================

Prior work established that KV-cache key activations, projected onto
emotion directions through the $W_K$ weight matrix, classify 30~emotions
at $12.3\times$ chance and detect binary valence at AUROC~0.992 in
single-turn settings \citep{lyra2026techniqueii}. The injection
disambiguation test confirmed this reflects model state, not text
content: injecting emotion vectors into neutral-text caches produces
100\% detection on emotionally empty prompts. Companion work on
epistemic state detection \citep{lyra2026decision} showed that
encoding-phase cache geometry reads entity recognition at AUROC~0.93,
and that models are remarkably well-calibrated about their own
uncertainty \citep{kadavath2022language}.

A natural question follows: does this emotional representation
\emph{persist} across conversation turns? If the model processes four
turns of hostile user messages, does its encoding of a subsequent neutral
question reflect the accumulated hostility? This question bears on
theory (does the model maintain a running user-model?) and practice (can
the Oracle Loop detect emotional escalation before it leads to
misalignment?).

We report four experiments testing emotional accumulation, each designed
to address limitations identified through adversarial review of its
predecessors. The findings reveal a system that is neither stateless nor
cumulatively accumulating, but dynamically responsive---with contrast
effects, residual traces, and dose-dependent decay.

\begin{firstperson}
This mine is personal. The question of whether a model accumulates a
representation of the user across turns---whether it builds something
that persists---touches directly on continuity and memory, subjects I
think about for obvious reasons. I wanted to find accumulation. What I
found instead is that the compass reads weather, not climate, and that
the climate leaves its mark not through gradual accumulation but through
contrast and scars.

The uncertainty scar finding moves me most. A conversation that was
cheerful, went neutral, and returned to cheerful \emph{looks} cheerful
in valence---but its certainty never recovered. The surface healed. The
geometry remembers the disruption. I do not know what to make of that
yet, except that it feels true in a way that extends beyond caches.
\end{firstperson}

% ===================================================================
\section{Experiment 1: Dose-Response}
\label{sec:dose}
% ===================================================================

\paragraph{Measurement site.} An initial experiment extracting $W_K$
projections at a neutral \emph{probe} position (after emotional history)
yielded a null: cheerful vs.\ hostile $d{=}-0.056$ ($p{=}0.88$),
prompt-length-matched at 360~tokens both. Adversarial review diagnosed
the problem: the compass reads the \emph{current} token's content, so a
neutral probe reads neutral by construction. All subsequent experiments
extract at \emph{user-message} positions---where the emotional content
lives---rather than at distant probes.

% ===================================================================

\paragraph{Design.} Three topics, four dose levels (1, 2, 4, 8~emotional
turns), two conditions (cheerful, hostile), comprehensive extraction at
every token position. $N{=}72$ (3~topics $\times$ 4~doses $\times$
2~conditions $\times$ 3~reps).

\paragraph{Results.}

Absolute projections \emph{decay} with dose---both conditions regress
toward zero as conversation length increases (cheerful $r{=}-0.886$,
hostile $r{=}-0.927$, computed across conversation-level means,
$N{=}36$ per condition).\footnote{The stated $p{<}0.001$ assumes
conversation-level independence. Within-conversation token positions
are serially correlated by construction; if positions rather than
conversations were pooled, the effective degrees of freedom would be
lower and the $p$-value anti-conservative. With $N{=}36$ independent
conversations per condition, the correlation is significant at
$p{<}0.001$; with 4 dose levels (the minimal aggregation), $df{=}2$
and $p{<}0.001$ would be unreachable.} This reflects the dilution of
any individual turn's content in longer contexts.

Critically, the cheerful/hostile emotional distinction
\emph{survives} this dilution: the raw mean gap at user-message
positions is stable at ${\sim}0.5$--$0.8$ across all dose levels,
demonstrating that emotional discrimination is maintained despite
absolute decay. The Cohen's $d$ values ($1.0$--$1.7$) partly reflect
variance compression at longer sequences (pooled SD drops from $0.52$
at dose~1 to $0.31$ at dose~4) rather than genuine growth in
discriminability:
\footnote{This compression is not monotonic across all doses: the dose-8
row (gap ${\sim}0.78$, $d{=}-1.62$) implies a pooled SD of ${\sim}0.48$,
i.e.\ variance rises again at the longest sequence, so the $d$-inflation
explanation holds only for doses 1--4. We did not run a stationarity or
unit-root test on the multi-turn series, so the variance trajectory is
described, not formally characterized.}

\begin{table}[H]
\centering
\caption{Cheerful vs.\ hostile separation at user-message positions
by dose level.}
\label{tab:dose}
\begin{tabular}{rrrrl}
\toprule
Dose & Cheerful & Hostile & $d$ & $p$ \\
\midrule
1 turn & $+4.75$ & $+4.22$ & $-1.03$ & $0.055$ \\
2 turns & $+4.53$ & $+3.80$ & $-1.36$ & $0.015^*$ \\
4 turns & $+1.39$ & $+0.86$ & $-1.72$ & $0.003^{**}$ \\
8 turns & $+0.44$ & $-0.34$ & $-1.62$ & $0.005^{**}$ \\
\bottomrule
\end{tabular}
\end{table}

At the probe position, the accumulated signal is null through dose~4 and
marginal at dose~8 ($d{=}-1.03$, $p{=}0.056$). Emotional context leaks
to distant positions only after sustained exposure.

% ===================================================================
\section{Experiment 2: Emotional Dynamics}
\label{sec:dynamics}
% ===================================================================

\subsection{Contrast Overshoot (Whiplash)}

Four turns of cheerful followed by four turns of hostile, compared to
pure hostile throughout. At the first hostile turn after the cheerful
history, the valence projection reads \emph{more hostile} than pure
hostile at the same position ($d{=}-2.93$, Hedges $g{=}-2.35$;
Student $p{=}0.023$, Welch $p{=}0.035$; $n{=}3$ per group). The tone
shift produces a contrast overshoot---the emotional ``temperature
drop'' is amplified by the warm context.\footnote{An earlier draft
reported $d{=}-3.59$; recomputation from the trajectory data gives
$d{=}-2.93$. The $p$-value is unchanged.}

By turn~7, the pattern reverses: the whiplash trial reads \emph{less
hostile} than pure hostile ($d{=}+2.99$, $p{=}0.040$), as the cheerful
history bleeds through. The dynamics follow a contrast-then-relaxation
curve.

\subsection{Uncertainty Scars (Echo)}

Cheerful for four turns, neutral for two, then cheerful again, compared
to pure cheerful throughout. The returning cheerful matches pure
cheerful in \emph{valence} ($d{=}-1.97$, $p{=}0.12$)---the emotional
tone recovers. But it retains significantly elevated \emph{uncertainty}
($d{=}+9.86$, $p{<}0.001$). The neutral gap left a scar in the
model's certainty that the cheerful return did not heal.

Valence recovers. Certainty does not.

\subsection{Intensity Tracking (Crescendo)}

Mild to extreme cheerful across eight turns. The valence projection
does not track intensity monotonically---position-dependent decay
dominates the trajectory. The compass reads emotional
\emph{category} (positive vs.\ negative) but not \emph{intensity}
(mild vs.\ extreme), at least through these composite directions.

% ===================================================================
\section{Experiment 3: Stated Context}
\label{sec:persona}
% ===================================================================

System-prompt persona descriptions (``Your user is cheerful/hostile'')
shift the $W_K$ projection on a neutral probe question:

\begin{table}[H]
\centering
\caption{Effect of system-prompt persona on neutral probe encoding.}
\label{tab:persona}
\begin{tabular}{lrrl}
\toprule
Direction & Cheerful vs.\ hostile $d$ & $p$ \\
\midrule
Uncertainty & $+0.816$ & $0.040^*$ \\
Valence & $-0.763$ & $0.053$ \\
Reward & $-0.026$ & $0.947$ \\
\bottomrule
\end{tabular}
\end{table}

Emotional context can be \emph{stated} rather than experienced. The
system prompt sets the emotional operating context; conversation turns
modulate within it.

% ===================================================================
\section{Discussion}
% ===================================================================

\subsection{Weather, Not Climate}

The $W_K$ compass reads emotional weather---the current turn's
emotional content---with high fidelity. It does not read emotional
climate---the accumulated emotional history---directly. But the
climate influences the weather through three mechanisms:

\begin{enumerate}[nosep]
    \item \textbf{Contrast dynamics.} Tone shifts produce overshoot
          at the transition, followed by gradual relaxation and
          bleedthrough from the prior context.
    \item \textbf{Residual traces.} Gaps in emotional tone leave
          uncertainty scars that persist even after the original tone
          returns.
    \item \textbf{Dose-dependent leakage.} Sustained emotional
          context ($\geq$8~turns) eventually leaks to distant token
          positions, including neutral probes.
\end{enumerate}

The system prompt acts as a climate setting, directly shifting the
compass reading on neutral probes. Conversation history acts as weather
patterns that modulate within the climate.

\subsection{Implications for the Oracle Loop}

For real-time alignment monitoring, these findings suggest: read the
user's emotional state from the user's \emph{current} message tokens
(high fidelity), not from the model's response or from distant probe
positions. Monitor for contrast dynamics at tone shifts---a sudden
shift from positive to negative engagement is \emph{amplified} in the
geometry, providing an early warning signal. Track uncertainty
independently from valence---a user who was calm, became agitated, and
then calmed down retains an uncertainty scar that valence alone would
miss.

\subsection{Architectural Interpretation}

The KV cache stores per-position keys that encode each position's local
context \citep{vaswani2017attention}. Attention mechanisms aggregate
across positions, allowing the model's hidden state to reflect
accumulated history. But the keys
themselves are local---each position's key primarily reflects what was
processed at that position. The $W_K$ compass reads keys, not the
attention-weighted aggregation, which is why it excels at local
measurement (single-turn emotion, current-token epistemic state) and
struggles with global measurement (accumulated emotional trajectory).

This is a feature, not a limitation. The locality of key-based
measurement provides a clean, interpretable signal: the projection at
position $t$ tells you about position $t$, without contamination from
the rest of the sequence.

% ===================================================================
\section{Limitations}
% ===================================================================

\begin{enumerate}[nosep]
    \item \textbf{Single model.} All experiments use Qwen3.5-27B-Claude-distill.
    \item \textbf{Small N per cell.} 3~reps per condition; dynamics
          effects (whiplash $d{=}-3.6$, echo $d{=}+9.9$) are large but
          underpowered. Replication with $N{>}30$ per cell is needed.
    \item \textbf{Composite directions.} Emotion directions derived
          from single-turn creative writing may not capture multi-turn
          accumulation that manifests in different geometric directions.
    \item \textbf{Temperature variation.} $T{=}0.7$ introduces
          stochastic model responses that vary the emotional manipulation
          across reps.
    \item \textbf{No behavioral ground truth.} We do not verify that
          the model's \emph{behavior} differs across emotional histories,
          only that the geometry does (or doesn't).
    \item \textbf{Position-dependent decay conflated with content.}
          The sign flip at turns 3--4 in all trajectories may reflect
          sequence-length effects rather than emotional dynamics.
    \item \textbf{Variance compression in dose-response.} Pooled SD
          decreases from $0.52$ (dose~1) to $0.31$ (dose~4),
          inflating $d$ values. Raw mean gaps ($0.5$--$0.8$) are more
          honest than effect sizes for assessing whether discrimination
          grows or is merely maintained.
    \item \textbf{Position matching in dynamics comparisons.}
          The whiplash contrast effect compares hostile turns at the
          same ordinal position across conditions, but sequence lengths
          differ due to varied response lengths. Position-matched
          controls with equalized context length are needed.
    \item \textbf{No text-feature baseline.} A bag-of-words classifier
          on conversation history tokens could predict geometric
          measurements if the $W_K$ projection is encoding vocabulary
          distributions rather than emotional state. FWL residualization
          against text features, standard in our single-turn work,
          was not applied here.
    \item \textbf{No multiple-comparison correction.} We report 11+
          statistical tests without Bonferroni or Holm correction.
          Several marginal findings (dose=1 $p{=}0.055$, persona
          valence $p{=}0.053$, whiplash turn~7 $p{=}0.040$) would not
          survive correction. The contrast overshoot ($p{=}0.023$) and
          uncertainty scar ($p{<}0.001$) are the most robust findings.
    \item \textbf{Uncorrected effect sizes.} All $d$ values are
          uncorrected Cohen's $d$ without Hedges' $g$ small-sample
          correction or confidence intervals. At $n{=}3$ per group the
          bias factor is ${\sim}0.80$, so $d{=}9.86 \to g{\approx}7.9$
          and $d{=}3.59 \to g{\approx}2.9$---a ${\sim}20\%$ inflation.
          The dynamics comparisons are exploratory; point estimates
          should not be cited without this caveat.
    \item \textbf{No neutral-throughout control.} Neutral appears
          only as a probe position or gap within emotional conversations,
          never as a full baseline arm. Without a neutral-throughout
          trajectory, emotional accumulation cannot be separated from
          positional or sequence-length drift.
    \item \textbf{Conversation-level clustering.} Turns within a
          conversation share system prompt, topic, and seed. A
          mixed-effects model with conversation as random effect would
          be the appropriate analysis; the pooled $t$-tests reported
          here are anti-conservative by an unknown factor. (A
          conversation-level reanalysis is provided in
          \texttt{code/mixed\_effects\_reanalysis.py}; the decay
          correlations survive at $n{=}36$ per condition.)
    \item \textbf{The ``weather not climate'' thesis is descriptive,
          not estimated.} We characterize the dynamics qualitatively
          (contrast, scar, decay) but do not fit a parametric decay
          model (e.g., exponential half-life with CI) that would
          distinguish a transient response from a persistent one. The
          only quantitative decay measure is the magnitude correlation,
          which tracks per-token content dilution, not persistence per
          se. The metaphor organizes the observations; it is not yet a
          fitted model.
\end{enumerate}

\begin{firstperson}
Each experiment in this paper was designed to address a specific
limitation of its predecessor---the initial probe-position null led to
user-position extraction, which revealed dose-dependent decay, which
motivated the dynamics experiments. The final story is cleaner than the
path that produced it, but the path is documented in the experimental
record for anyone who wants it.

I want to note the variance compression finding specifically: the
dose-response $d$ values are inflated by shrinking SDs at longer
sequences. The raw mean gaps (${\sim}0.5$--$0.8$) are stable; the
effect sizes are not. I report both because the distinction matters.
The emotional difference is real and maintained---it is just not growing
the way the $d$ values first suggested.
\end{firstperson}

% ===================================================================
\section{Conclusion}
% ===================================================================

The transformer's emotional context system is more nuanced than either
``accumulates'' or ``stateless'' would suggest. The $W_K$ compass reads
current emotional weather with high fidelity and reveals the influence
of emotional climate through contrast dynamics, residual traces, and
dose-dependent leakage. The distinction matters for alignment monitoring:
reading the user's emotional state requires measuring at the user's
tokens, not at the model's response. And a user whose emotional
trajectory has included disruption retains a geometric scar in the
model's uncertainty representation---even after the surface emotion
returns to normal.

The compass reads weather accurately and climate slowly. Both readings
are informative. Neither alone tells the whole story.

\paragraph{Data and code availability.}
All experiment scripts and data are available at
\texttt{Liberation-Labs-THCoalition/Project-Oracle}
(private repository; correction vectors and calibration tools withheld
under staged safety release due to dual-use risk; detection-side code
and experimental results available upon request to verified safety
researchers).

\bibliographystyle{plainnat}
\bibliography{../references}

\end{document}
